% Generated draft: requires mathematical and source review.
\documentclass[11pt]{article}
\usepackage[margin=25mm]{geometry}
\usepackage{fontspec}
\setmainfont{Latin Modern Roman}
\usepackage{amsmath,amssymb,mathtools,unicode-math}
\setmathfont{Latin Modern Math}
\usepackage{enumitem,needspace,xurl,hyperref}
\setlength{\parindent}{0pt}
\setlength{\parskip}{5pt}
\setlength{\emergencystretch}{3em}
\newcommand{\N}{\mathbb{N}}
\newcommand{\Z}{\mathbb{Z}}
\newcommand{\Q}{\mathbb{Q}}
\newcommand{\R}{\mathbb{R}}
\newcommand{\C}{\mathbb{C}}
\newcommand{\F}{\mathbb{F}}
\newcommand{\A}{\mathbb{A}}
\newcommand{\PP}{\mathbb{P}}
\newcommand{\E}{\mathbb{E}}
\newcommand{\eps}{\varepsilon}
\newcommand{\dd}{\,\mathrm{d}}
\begin{document}
\section*{UnsolvedMath ID: 30001263}
\textbf{Source problem code:} catalogue entry
\subsection{Definitions and mathematical statement}
Let \(G\) be a torsion-free discrete group, meaning that its identity is its only element of finite order, and let \(R\) be a regular ring in the sense intended by the supplied source. The precise convention for regularity is not specified in the record. Write \(RG\) for the group ring, whose elements are finite sums \(\sum_{g\in G}r_g g\), with multiplication induced by the multiplication in \(R\) and the group law in \(G\). Let \(BG\) denote a classifying space for \(G\). For a spectrum \(\mathbf E\), write \(H_n(BG;\mathbf E)\) for generalized homology with constant spectrum coefficients. Let \(\mathbf K_R\) denote the nonconnective algebraic \(K\)-theory spectrum of \(R\), and write \(K_n(RG)\) for the corresponding algebraic \(K\)-groups. The notation \(\mathbf L_R^{\langle-\infty\rangle}\) and \(L_n^{\langle-\infty\rangle}(RG)\) denotes the ultimate lower-decorated algebraic \(L\)-theory spectrum and groups intended by the source; the necessary involution and precise \(L\)-theory conventions are not supplied.

The displayed question in the supplied clean statement is whether, for every torsion-free discrete group \(G\), every regular ring \(R\), and every integer \(n\), both indicated assembly homomorphisms are isomorphisms:
\[
 A^K_{G,R,n}\colon H_n(BG;\mathbf K_R)\longrightarrow K_n(RG),
 \qquad
 A^L_{G,R,n}\colon H_n(BG;\mathbf L_R^{\langle-\infty\rangle})\longrightarrow L_n^{\langle-\infty\rangle}(RG).
\]
Here ``isomorphisms'' expresses the source's requirement that the maps be bijective.

\emph{Unresolved formulation qualifications.} The record does not specify whether rings are assumed unital, what regularity means, or what involution and associated conventions define the displayed \(L\)-theory terms and map. Consequently, the record does not determine a fully specified universal assertion. It also does not explain the hypotheses or comparison needed to identify these maps from \(BG\) with a formulation of the Farrell--Jones conjectures using assembly relative to the family of virtually cyclic subgroups; that identification is not asserted here.

The original statement has a broader scope: it asks whether the Kaplansky, negative-\(K\) and Whitehead, Novikov, Borel, Serre, Poincar\'e duality, and Farrell--Jones conjectures hold for unspecified ``stated classes'' of discrete groups and rings, highlighting the displayed assembly question in particular. The record does not give the individual formulations or the classes and hypotheses for those additional questions. They are retained as unresolved scope qualifications and cannot be replaced by the displayed assembly question without further source clarification.
\subsection{Short English statement}
The clean statement asks whether both displayed algebraic \(K\)- and \(L\)-theory assembly maps are isomorphisms in every integer degree for every torsion-free discrete group and every regular ring. Ring regularity and the data needed for \(L\)-theory remain unspecified. The original record also bundles several other conjectures without supplying their precise formulations. This is a statement-only draft requiring human mathematical review.
\subsection{Sources}
\begin{itemize}
\item Supplied local UnsolvedMath registry record, permanent ID \(30001263\), exact source problem code \texttt{catalogue entry}, titled ``Farrell--Jones Assembly Maps for Torsion-Free Groups.'' The record belongs to the Oberwolfach Reports Open Problems collection and gives proposed year 2009; no individual proposer is supplied.
\item \emph{Manifold Perspectives} (2009), Oberwolfach Reports, pp.~1487--1546, DOI 10.4171/owr/2009/27. Supplied source report and literature-source link labelled ``original OWR report'': \url{https://doi.org/10.4171/owr/2009/27}. The supplied background describes Oberwolfach Reports as documenting mathematical workshops and their open-problem discussions. No item-specific page or passage is supplied.
\item Supplied literature-source link labelled ``Lueck report on the conjectures'': \url{https://ems.press/content/serial-article-files/46228}. Further bibliographic details are not supplied.
\item The exact supplied literature-check date is \texttt{2026-08-21}. The supplied dated curation labels the entry open and assesses it as a bundle of conjectures, stating that Kaplansky and the general Novikov, Borel, and Farrell--Jones statements remain unresolved. This is supplied dated curation, not an independently verified literature assessment; the background explicitly requests independent verification of time-sensitive status and cited claims.
\item The registry labels the statement \texttt{corrected\_verified} and reports that it was checked against the official Oberwolfach report, preserving the stated hypotheses, notation, and open target. These are supplied verification claims; neither linked source was independently consulted in preparing this draft.
\end{itemize}
\end{document}
